\newpage
\section{Acta Número 22}
\label{sec:acta22}

\noindent \textbf{Fecha de la sesión:}\par
Lunes, 18 de Mayo de 2026

\vspace*{0.5cm}
\noindent \textbf{Datos de la Sesión:}
\vspace{-0.2cm}
\begin{itemize}[noitemsep]
    \item \textbf{Modalidad:} Virtual - Google Meet.
    \item \textbf{Hora de Inicio:} 18:00
    \item \textbf{Hora de Fin:} 20:00
\end{itemize}

\vspace*{0.5cm}
\noindent \textbf{Orden del Día:}
\vspace{-0.2cm}
\begin{itemize}[noitemsep]
    \item Control de Asistencia.
    \item Ceremonia de Inicio de Aseguramiento de Calidad (QA): Planificación y validación técnica.
    \item Ceremonia de Gestión de Errores y Correcciones: Protocolo de remediación y flujo de trabajo.
    \item Evaluación del estado de desarrollo y consenso sobre la prórroga para informes de calidad.
    \item Firmas y Aprobación de los Presentes.
\end{itemize}

\vspace*{0.5cm}
\noindent \textbf{Socios Fundadores Presentes:}
\vspace{-0.2cm}
\begin{enumerate}[noitemsep]
    \item Pardo Pozo Juan Diego (Jefe de Proyecto).
    \item Arnez Jaimes Mauricio.
    \item Quispe Flores Camila Belen.
    \item Ariñez Bautista Jim Gabriel.
    \item Medina Rodríguez Mateo Jhoustin.
    \item Molina Maldonado Tomas Manuel.
    \item Gutierrez Guzmán Angheline.
\end{enumerate}

\newpage
\subsection{Desarrollo del Acta}

\noindent \textbf{Control de Asistencia del Team}\par
Se procedió a la toma de asistencia a la hora acordada, corroborando la presencia de los 7 socios fundadores de la grupo-empresa Byte Busters S.R.L. Tras verificar el quórum legal y la disposición de todos los integrantes, se dio inicio formal a la sesión de planificación de calidad.

\vspace*{0.5cm}
\noindent \textbf{Inicio de Aseguramiento de Calidad (QA): Rigurosidad y Criterios}\par
El equipo inauguró formalmente la fase de Aseguramiento de Calidad para el presente ciclo. Durante esta ceremonia, se definió la hoja de ruta para la ejecución de pruebas unitarias e integrales, asegurando que cada componente del sistema responda a los exigentes criterios de aceptación establecidos. Se enfatizó la importancia de la validación técnica exhaustiva y se acordó un protocolo ágil para la aplicación de \textit{Hotfixes} (correcciones de errores críticos) en tiempo real, con el fin de no comprometer la estabilidad de las ramas principales de desarrollo.

\vspace*{0.5cm}
\noindent \textbf{Gestión de Errores y Correcciones: Flujo de Remediación}\par
Se procedió a socializar el protocolo de gestión de incidencias, el cual garantiza una trazabilidad total desde la detección hasta la resolución. El flujo acordado comprende desde la reproducción y el triaje (priorización) de errores, pasando por un análisis profundo de causa raíz (\textit{RCA - Root Cause Analysis}), hasta la implementación de la solución y las pruebas de regresión local. Finalmente, se estableció que toda corrección debe ser re-entregada al equipo de QA mediante un \textit{push} formal para su validación definitiva, cerrando así el ciclo de mejora continua.

\vspace*{0.5cm}
\noindent \textbf{Estatus del Desarrollo y Acuerdo de Prórroga Técnica}\par
En un ejercicio de transparencia y honestidad técnica, se evaluó el progreso individual de las tareas asignadas. Al detectarse que algunos módulos de desarrollo aún requieren ajustes finos y que los informes de calidad demandan una redacción detallada, el equipo decidió por unanimidad conceder una **semana adicional de desarrollo**. Este periodo será utilizado estratégicamente para culminar las funcionalidades pendientes y perfeccionar la documentación técnica y los reportes de QA, garantizando que el entregable final cumpla con los estándares de excelencia de Byte Busters S.R.L.

\vspace*{0.5cm}
\noindent \textbf{Aprobación Oficial}\par
Tras haber resuelto todas las dudas sobre la semana de QA y habiendo llegado a un acuerdo mutuo sobre la extensión del plazo de desarrollo, todos los miembros fundadores expresan su conformidad con los hitos planteados. Sin más asuntos que tratar, se procede con la firma de los presentes.

\newpage
\noindent \textbf{Firmas y Aprobación de los Presentes}
\begin{center}
    \noindent
    \setlength{\tabcolsep}{0pt}
    \begin{tabular}{ccc}
        \espacioFirma{assets/img/firmas/mauricio.jpg}{Mauricio Jaimes Arnez}{13751221} & 
        \espacioFirma{assets/img/firmas/camila.jpg}{Camila Belen Quispe Flores}{13097244} &
        \espacioFirma{assets/img/firmas/jim.jpg}{Jim Gabriel Ariñez Bautista}{9517642} \\[1.0cm]
        \espacioFirma{assets/img/firmas/mateo.jpg}{Mateo Medina Rodríguez}{9453809} &
        \espacioFirma{assets/img/firmas/tomas.jpg}{Tomas Molina Maldonado}{8051701} &
        \espacioFirma{assets/img/firmas/angheline.jpg}{Angheline Gutierrez Guzmán}{13751815} \\[1.0cm]
    \end{tabular}
    \vspace{0.2cm} 
    \begin{tabular}{cc}
        \espacioFirma{assets/img/firmas/diego.jpg}{Juan Diego Pardo Pozo}{12747783}
    \end{tabular}
\end{center}

\vspace{0.5cm}
\noindent \textbf{Fecha de última revisión:} Lunes, 18 de Mayo de 2026